%%%%%%%%%%%%%%%%%%%%%%%%%%%%%%%%%%%%%%%%%%%%%%%%%%%%%%%%%%%%%%%%%%%%%%%%%%%%%%%
\chapter{Introduction} %%%%%%%%%%%%%%%%%%%%%%%%%%%%%%%%%%%%%%%%%%%%%%%%%%%%%%%%

\section{Motivation}

Rowing is a technically demanding sport. Good performance does not only depend on strength or endurance, but also on how well the athlete coordinates the movement during each stroke. A rowing stroke is usually described by the phases catch, drive, finish, and recovery. In these phases the athlete, the sliding seat, the oars, and the boat form one coupled mechanical system. Small changes in timing, force application, body position, or boat balance can therefore influence the efficiency of the stroke. Reviews of rowing technique show that coaches and researchers have long tried to connect stroke mechanics with boat velocity and performance, but also that it is difficult to define one single ideal technique that works for every athlete and setup \cite{soper_towards_2004}.

In everyday training, rowing technique is often evaluated by a coach watching the athlete, by video, or by commercial boat instruments. These methods are useful, but they have limits. Visual feedback is subjective and depends on the experience of the coach. Video can be detailed, but it is time-consuming to analyze and does not always give numerical values. Commercial systems can measure useful quantities such as boat speed, stroke rate, oar angle, or force, but they can be expensive, boat-specific, or complex to install. This creates a practical gap between laboratory-grade motion analysis and simple training feedback on the water.

Inertial measurement units (IMUs) are attractive for this problem because they are small, low-cost, and can measure motion directly on the athlete, the equipment, or the boat. IMUs are already used in many sports and movement-analysis systems because they can record acceleration and angular velocity without cameras or a fixed laboratory setup \cite{arlotti_benefits_2022}. In rowing, IMUs have been used to measure boat motion, oar motion, handle motion, and athlete movement \cite{llosa_remote_2009,junker_sensor-based_2015,hohmuth_wireless_2023,rodriguez_mendoza_wearable_2024}. This makes them a strong candidate for a low-cost rowing prototype that can be used outside the lab.

However, IMUs also have a central weakness. They do not directly measure position. Position must be estimated by integrating acceleration, and small sensor errors, bias, orientation errors, and noise grow over time during this integration. This is a well-known problem in inertial sensing: inertial estimates can be accurate over short time intervals, but they drift over longer time intervals unless they are corrected with models, constraints, or additional sensors \cite{rodriguez_mendoza_wearable_2024}. For this reason, many advanced systems use sensor fusion, for example by combining inertial sensors with GNSS, UWB, external reference systems, or biomechanical models \cite{cloud_adaptive_2019,rodriguez_mendoza_wearable_2024,ruffaldi_sensor_2015}.

This thesis follows a simpler and more robust idea. Instead of trying to reconstruct the exact absolute position of the rowing seat, the prototype focuses on relative and stroke-level metrics. The main target is not to say: ``the seat is exactly at this position in meters''. The target is to measure signals that are useful for rowing feedback, such as forward acceleration, stroke timing, movement intensity, pitch and yaw stability, and the difference between a seat-mounted sensor and a boat-mounted reference sensor. In the two-sensor case, the core idea is:

\begin{equation}
  \text{relative motion} = \text{SEAT signal} - \text{BOAT signal}.
\end{equation}

This is important because a seat-mounted IMU on its own measures both the athlete's motion and the motion of the boat under the athlete. A boat-mounted IMU can provide a reference for the boat motion. Subtracting the boat signal from the seat signal can therefore help isolate movement that is more related to the athlete and sliding seat. This principle is close to the seat-position approach by Gravenhorst et al., who used one sensor on the sliding seat and one sensor on the boat shell to calculate seat acceleration relative to the boat before further processing \cite{gravenhorst_self-aligning_2012}. The difference in this thesis is that the prototype does not need to solve the full problem of accurate double-integrated seat position. It can avoid part of the drift problem by using relative metrics and per-stroke signal features.

\section{Problem Statement}

The problem addressed in this thesis is the development of a low-cost IMU-based sensing prototype for rowing. The system should capture rowing-related movement data in a practical way and should support later analysis of stroke-level metrics. It should also be suitable for use in a realistic rowing environment, where water, vibration, mounting constraints, battery life, and device size matter.

This problem has both a sensing side and a hardware-integration side. On the sensing side, the system must record data at a sufficient rate, store or transmit the data reliably, and keep the data format simple enough for later analysis. On the hardware side, the electronics must be protected and mounted in a way that does not disturb the athlete or change the rowing movement. Waterproof and robust enclosure design is therefore not only a mechanical detail. It can directly influence data quality, because unstable mounting, water ingress, or loose components can create artifacts in the sensor signal. Work on waterproof sensing enclosures shows that sealing, material choice, assembly, and iterative prototype testing are important parts of making measurement devices usable outside controlled laboratory conditions \cite{lindqvist_design_2015,huayna_fernandez_enclosure_2025}.

The rowing environment makes the problem harder than many indoor motion-analysis tasks. The boat is not a fixed reference frame. It accelerates and decelerates during the stroke. It also moves due to water, wind, steering, and the athlete's body mass. A sensor on the athlete or seat therefore observes a mixture of several motions. If the goal were exact absolute seat position, the system would need careful calibration, synchronization, drift correction, and validation against a reference system. For a bachelor-level prototype, this would be a high-risk target.

For this reason, the thesis uses a more modest but still useful goal. The prototype is designed to collect reliable IMU data and derive relative metrics that can support rowing feedback. These metrics are easier to obtain than absolute position, because they can be based on signal shape, peaks, timing, standard deviation, root mean square values, and relative differences between sensors. Such metrics still describe meaningful parts of the rowing motion, especially when they are compared between strokes, between sessions, or between a seat sensor and a boat reference sensor.

\section{Approach}

The developed prototype is based on a small embedded sensing unit with an IMU, SD card logging, USB serial output, and Bluetooth Low Energy (BLE) streaming. The current embedded software records acceleration and angular velocity data at approximately 100 Hz. The data are written as comma-separated values and include a device identifier, sequence number, timestamp, acceleration axes, and gyroscope axes. The software supports a device role such as \texttt{SEAT}, and the same code can be adapted for a second unit such as \texttt{BOAT}.

The software side contains three main tools. First, a live serial or BLE reader can receive the CSV stream and calculate simple rolling metrics. Second, a live web dashboard can display the current signal, record a training segment, and mark good or bad strokes. Third, an offline analysis script can load a recorded CSV file, check the sample timing, detect sequence gaps, calculate summary statistics, estimate stroke timing, and generate an HTML report. In the two-IMU mode, the live dashboard can connect to a seat unit and a boat unit and compute relative signals from the difference between both streams.

The main design choice is to treat the IMU data as a source for robust training indicators rather than as a perfect position sensor. The analysis therefore uses the mounted X axis as a forward/backward candidate axis, smooths and centers the signal for stroke detection, and calculates stroke rate and movement proxies. A relative velocity proxy can be estimated by integrating a smoothed acceleration segment with simple drift correction, but it is clearly treated as a proxy and not as an absolute speed in meters per second. A power-transfer proxy can be calculated from the product of forward acceleration and relative speed, but it is also not interpreted as real mechanical power in watts.

This approach is intentionally conservative. It accepts the known limits of low-cost IMUs and avoids overclaiming. At the same time, it keeps the system useful for a first prototype: the system can still show whether the movement is periodic, whether strokes are detected consistently, whether the boat reference signal differs from the seat signal, and whether relative movement patterns are stable across strokes.

\section{Contributions}

The thesis contributes a practical prototype and analysis pipeline for low-cost rowing motion sensing. The main contributions are:

\begin{itemize}
  \item an embedded IMU logger for rowing-related acceleration and angular velocity data;
  \item a simple CSV-based data format with timestamps and sequence numbers for later analysis;
  \item USB serial and BLE streaming for live monitoring;
  \item an offline analysis workflow for sampling checks, signal statistics, stroke detection, and report generation;
  \item a two-IMU relative measurement concept using a seat-mounted and boat-mounted sensor;
  \item a metric strategy that focuses on relative and stroke-level indicators instead of exact absolute seat position;
  \item a hardware-integration perspective that treats enclosure design, mounting, and robustness as part of measurement quality.
\end{itemize}

The expected result is not a final commercial rowing coach. It is a research prototype that can be evaluated, improved, and used to collect first data. The value of the work lies in connecting requirements, hardware design, embedded sensing, data processing, and rowing-specific interpretation in one system.

\section{Thesis Structure}

The rest of this thesis is structured as follows. Chapter~2 reviews related work on rowing technique, IMU-based rowing systems, seat tracking, sensor fusion, general IMU limitations, and enclosure design. Chapter~3 describes the system requirements and design decisions. Chapter~4 explains the implementation of the embedded logger, data pipeline, live dashboard, and offline analysis. Chapter~5 presents the prototype case and experimental evaluation. Chapter~6 discusses the results, limitations, and future work. Chapter~7 concludes the thesis.


%%%%%%%%%%%%%%%%%%%%%%%%%%%%%%%%%%%%%%%%%%%%%%%%%%%%%%%%%%%%%%%%%%%%%%%%%%%%%%%
\chapter{Related Work} %%%%%%%%%%%%%%%%%%%%%%%%%%%%%%%%%%%%%%%%%%%%%%%%%%%%%%%%

\section{Rowing Technique and the Need for Feedback}

Rowing performance depends on a coordinated stroke cycle. In sculling and sweep rowing, the athlete must move the body, seat, oars, and boat in a controlled sequence. Soper and Hume describe rowing performance as a complex interaction of kinematics, kinetics, anthropometry, boat setup, and coaching practice \cite{soper_towards_2004}. They also point out that rowing biomechanics has investigated many aspects of technique, but that simple and general predictors of performance are difficult to define. This is important for the present thesis because it shows that a sensing system should not claim to measure ``perfect technique'' directly. It should instead provide objective data that can help interpret technique.

Hohmuth et al. also describe rowing as a sport in which technique, timing, force, and coordination strongly influence performance \cite{hohmuth_wireless_2023}. Their work lists typical technique-related variables such as stroke length, blade depth, catch and finish timing, smooth recovery, force application, seat displacement, and boat acceleration. This supports the idea that a useful rowing sensing prototype should not only measure one single number. It should allow several views of the stroke, such as timing, forward movement, stability, and periodic signal shape.

Feedback systems in rowing are especially valuable because the athlete cannot easily see the own movement while rowing. Coaches often use observation and video, but objective sensor data can help make repeated training sessions more comparable. Junker's doctoral work on sensor-based performance feedback in rowing shows the broader research direction of using sensors to create quantitative feedback for rowing technique \cite{junker_sensor-based_2015}. This thesis follows the same general direction, but with a smaller and cheaper prototype focused on IMU-based relative metrics.

\section{Wireless and IMU-Based Rowing Measurement Systems}

One early example of a wireless rowing sensing system is REMOTE by Llosa et al. \cite{llosa_remote_2009}. REMOTE uses wireless sensor network nodes with accelerometers to monitor oar and boat motion. The system was tested through laboratory calibration, ergometer tests, and on-boat tests. The work shows that accelerometer-based wireless systems can provide useful information such as boat motion, oar motion, stroke angles, boat balance, and stroke trajectory. It also shows a limitation that is highly relevant for this thesis: the authors attribute deviations in stroke trajectory to handmade packaging, unknown original position, and numerical integration errors. This is a clear example of how mechanical integration and inertial processing errors can influence the final motion estimate.

The same general rowing-sensing direction was later extended in several IMU-based systems. Gravenhorst et al. validated an IMU-based rowing oar angle measurement system and compared the IMU result with a reference system based on a potentiometer \cite{gravenhorst_validation_2013}. Oar-mounted IMUs are useful because oar angles are strongly connected to stroke length, timing, and rowing phases. However, oar-angle systems mainly describe the equipment motion. They do not directly isolate the movement of the rower's seat or body relative to the boat.

Avvenuti et al. present MARS, a multi-agent system for assessing rower coordination using motion-based stigmergy \cite{avvenuti_mars_2013}. Their system focuses on crew synchrony rather than exact biomechanical reconstruction. This is useful for the present thesis because it shows another practical strategy: rowing feedback can be based on robust patterns and timing indicators, not only on full 3D position reconstruction. A simple metric can still be valuable if it is stable, interpretable, and connected to the coach's view of rowing quality.

Hohmuth et al. developed WiRMS, a wireless rowing measurement system that combines several sensing modalities \cite{hohmuth_wireless_2023}. Their system includes handle units with IMUs, strain gauges, and pressure sensors, a boat kinematic measurement unit with acceleration and GPS, contactless seat displacement measurement, and synchronized EMG data. A proof-of-concept study with experienced athletes showed that the system can record movement and force variables during on-water rowing. WiRMS is much richer than the prototype in this thesis, but it also shows the cost of richer measurement: more sensors, more integration work, and more setup complexity. The present thesis therefore aims for a smaller system that is easier to build and evaluate.

\section{Seat Motion and Seat-to-Boat Relative Measurement}

The sliding seat is an important part of rowing technique because it is directly connected to leg drive and body movement. Gravenhorst et al. argue that seat movement is linked to the rower's leg movement and therefore to one of the most powerful parts of the rowing stroke \cite{gravenhorst_self-aligning_2012}. They also note that simple event-based devices can measure stroke rate but not the full shape of the seat motion.

The most important related work for this thesis is the self-aligning and drift-compensated rowing seat position measurement system by Gravenhorst et al. \cite{gravenhorst_self-aligning_2012}. Their system uses two miniaturized sensor units: one attached to the sliding seat and one attached to the boat shell. The seat sensor measures the superposition of boat acceleration and seat acceleration relative to the boat. The boat sensor measures boat acceleration. The relative seat acceleration can therefore be calculated from the difference between both sensors, after correcting for sensor misalignment. The authors then use drift compensation and double integration to estimate seat position.

This idea is directly relevant because it gives a strong technical reason for a two-IMU topology. A single sensor on the seat cannot easily separate athlete or seat movement from boat movement. A second sensor on the boat can act as a reference. However, the work also shows the difficulty of the full seat-position problem. To obtain position from acceleration, the signal must be integrated twice, and low-frequency drift must be corrected. Gravenhorst et al. use discrete correction events to support drift compensation \cite{gravenhorst_self-aligning_2012}. The present thesis takes inspiration from their relative measurement idea but avoids the stronger claim of accurate absolute seat position. It focuses on relative acceleration, stroke timing, and proxy metrics.

Gravenhorst et al. also developed SonicSeat, a contactless seat position tracker based on ultrasonic sound measurements \cite{gravenhorst_sonicseat_2012}. SonicSeat measures seat position directly with an ultrasonic distance sensor mounted in the boat. The system was used to derive metrics such as stroke rate, catch delay, maximum sliding velocity, and sliding seat displacement. The work is important because it shows that seat motion can be translated into coach-relevant metrics. It also shows that direct distance sensing can avoid part of the inertial double-integration problem. At the same time, ultrasonic sensing requires additional hardware and still needs a stable installation in the boat. The IMU-based prototype in this thesis is less direct, but smaller and more general.

\section{Sensor Fusion and Positioning for Rowing}

Some rowing systems try to estimate richer kinematic variables by fusing several sensor types. Cloud et al. use smartphone GPS and accelerometer data to estimate boat speed, distance, and distance per stroke \cite{cloud_adaptive_2019}. They compare a complementary filter and a Kalman filter against differential GPS reference data. Their results show that activity-specific sensor fusion can improve estimates from low-cost consumer sensors. At the same time, they report that distance-per-stroke accuracy remains limited compared with the needs of very detailed stroke-level analysis. This is important for the present thesis because it shows both sides of low-cost sensing: useful improvements are possible, but exact high-accuracy motion estimates remain difficult.

Rodriguez Mendoza and O'Keefe present a wearable multi-sensor positioning prototype for rowing technique evaluation \cite{rodriguez_mendoza_wearable_2024}. Their system combines an IMU, ultra-wideband ranging, and GNSS to estimate the position, velocity, and attitude of a wearable node during rowing. They compare several methods, including UWB, inertial navigation, rowing dead reckoning, GNSS-aided inertial navigation, and UWB-aided inertial navigation. The paper explicitly states that inertial navigation solutions degrade over time because acceleration and angular-rate errors are integrated into velocity, position, and attitude. The work therefore uses external aiding sources and a reference trajectory. This supports the decision in this thesis to avoid absolute position as the main output of a simple low-cost IMU prototype.

Ruffaldi et al. investigate sensor fusion for complex articulated body tracking in rowing \cite{ruffaldi_sensor_2015}. Their work represents a more complete body-tracking direction, using several IMUs and constraints to estimate articulated motion. This can provide richer information than a two-sensor seat/boat prototype, but it also requires more sensors, calibration, modeling, and validation. For a practical bachelor prototype, a smaller sensor set is more realistic. The trade-off is that the output metrics must be simpler and should be described as relative indicators.

\section{General Challenges of IMU-Based Motion Analysis}

The main technical challenge of IMU-based motion analysis is drift. Accelerometers and gyroscopes are useful because they provide high-rate motion data, but their errors accumulate when the data are integrated. Bias, noise, temperature effects, wrong initial orientation, and mounting errors can all affect the result. This is why many IMU systems use additional sensors, constraints, or event-based corrections. In rowing, the periodic structure of the stroke can help because peaks or repeated phases can be used for segmentation or reset logic \cite{cloud_adaptive_2019,rodriguez_mendoza_wearable_2024}.

Other sports show the same pattern. Yu et al. use an IMU-based system for football kick motion analysis and show a full processing chain including calibration, orientation estimation, and validation against camera-based data \cite{yu_motion_2022}. He et al. validate an IMU-based gait analysis system against optical motion capture, which is a common way to test whether wearable movement metrics are accurate enough \cite{he_accuracy_2024}. These works are not rowing systems, but they show that IMU metrics need careful validation and that accuracy claims should be tied to the measured variable.

Reviews of IMU use in sports and sports medicine also describe the practical benefits and limitations of wearable sensors. Arlotti et al. discuss IMU-based wearables as useful tools for sport and medical contexts, but the broader field still has to handle calibration, drift, placement, and interpretation issues \cite{arlotti_benefits_2022}. Presicci et al. review IMU-based systems in skiing and show that IMUs are widely used for sport-specific metrics, but that protocols and validation still matter strongly \cite{presicci_imu-based_2025}. These lessons transfer to rowing: a metric can be useful only if the mounting, processing, and interpretation are clear.

This thesis therefore uses a cautious interpretation of IMU data. It does not present integrated acceleration as exact seat position. It uses acceleration, angular velocity, signal variation, stroke events, and SEAT-minus-BOAT differences as relative indicators. This reduces the dependence on long-term integration and makes the prototype more suitable for early evaluation.

\section{Wearable Integration, Placement, and Enclosure Design}

Wearable sensing is not only an algorithmic problem. The sensor must be placed and mounted in a way that is comfortable, repeatable, and robust. Geissler et al. show an alternative wearable sensing approach using textile capacitive sensors embedded in garments \cite{geisler_embedding_2024}. Their work is not an IMU rowing system, but it highlights an important general point: sensing quality is connected to the physical integration of the sensor into a wearable object. If the sensor moves relative to the body or object, the measured signal changes.

This is also true for a rowing IMU prototype. A sensor on the seat or boat must be attached firmly enough that the signal represents the motion of the seat or boat, not vibration of the casing. At the same time, the device should be small and unobtrusive. It must survive water, splashes, mechanical stress, and handling before and after training. Waterproof enclosure design is therefore part of the measurement system. Lindqvist and Stroemberg Groendahl describe a waterproof probe enclosure development process with requirements, concept selection, O-ring sealing, material choices, and prototype testing \cite{lindqvist_design_2015}. Huayna Fernandez similarly focuses on enclosure development for a water-quality measuring tool and discusses the design of a housing for use in a wet environment \cite{huayna_fernandez_enclosure_2025}. These works support the thesis focus on hardware integration and casing design as more than cosmetic work.

The role of intelligent interpretation and feedback is also relevant. Sorysz discusses opportunities and challenges for artificial intelligence in wearables for physical therapy and sports training \cite{sorysz_artificial_2024}. For this thesis, this does not mean that a complex AI model is required. Instead, it shows that wearable sports systems should be designed with later feedback and interpretation in mind. A clean data format, stable mounting, and meaningful labels such as good or bad strokes can support future learning-based analysis.

\section{Research Gap}

The literature shows that rowing sensing is already a rich research field. There are systems for boat speed and acceleration, oar angle, handle force, pressure distribution, seat displacement, EMG, crew synchrony, and 3D handle positioning \cite{llosa_remote_2009,avvenuti_mars_2013,cloud_adaptive_2019,hohmuth_wireless_2023,rodriguez_mendoza_wearable_2024}. However, many of these systems are complex, expensive, intrusive, or focused on a different measurement target. Some focus mainly on oars or handles. Others use UWB, GNSS, EMG, multiple IMUs, or contactless distance sensors. These systems can provide rich data, but they increase setup effort and make the system harder to reproduce as a low-cost prototype.

There is also a technical gap between raw IMU signals and reliable rowing metrics. A low-cost IMU can measure acceleration and angular velocity well enough for many short-term patterns, but exact absolute position is hard because of drift and integration errors \cite{cloud_adaptive_2019,rodriguez_mendoza_wearable_2024}. This is especially important for the sliding seat, because a seat-mounted sensor measures both seat motion and boat motion. The two-sensor seat/boat concept by Gravenhorst et al. is therefore a strong starting point \cite{gravenhorst_self-aligning_2012}. The remaining opportunity is to build a simple prototype that uses this relative idea without depending on full double-integrated seat position.

This thesis addresses that opportunity. It develops a low-cost IMU-based prototype for rowing and focuses on relative, interpretable, and stroke-level metrics. The system is not designed to replace laboratory motion capture or high-end rowing instrumentation. Instead, it is designed as a practical step from requirements to prototype evaluation: record useful IMU data, protect the electronics, stream and store the signal, analyze stroke-level patterns, and use a boat reference sensor to reduce the ambiguity of single-sensor seat measurements.
